\scanpage{24}\mspage{9}
\uline{et pour différentielle invariante} $\pi=\mathrm{dx}/2\mathrm{y}$.

De plus, sous ces hypothèses, la donnée d'une forme différentielle invariante suffit à rigidifier une courbe.

4.\qquad \uline{Cas où} $2$ \uline{est nilpotent et} $\mathrm{c}_{4}$ \uline{inversible}.

Lorsque $2=0$, la forme différentielle ne suffit plus jamais à rigidifier une courbe de genre un. Lorsque $\mathrm{c}_{4}$ est inversible et $2$ nilpotent, donc $\mathrm{a}_{1}$ inversible, en changeant le choix de $\pi$, $\mathrm{x}$, $\mathrm{y}$ on peut trouver une équation
\[
\mathrm{y}^{2}+\mathrm{xy}=\mathrm{x}^{3}+\mathrm{a}_{2}\,\mathrm{x}^{2}+\mathrm{a}_{6}\qquad
\pi=\handwritten{\frac}{\mathrm{dx}}{\mathrm{x}+2\mathrm{y}}\qquad.
\tag*{(4.1)}
\]

Les autres choix de $\mathrm{x}$ et $\mathrm{y}$ donnant lieu à une telle équation s'obtiennent en faisant, dans (1.2)
\[
\mathrm{r}=\mathrm{t}=0\qquad\mathrm{u}=1+2\mathrm{s}
\tag*{(4.2)}
\]
et on aura
\[
\begin{aligned}[t]
\mathrm{u}^{2}\,\mathrm{a}'_{2}&=\mathrm{a}_{2}-\mathrm{s}-\mathrm{s}^{2}\\
\mathrm{u}^{6}\,\mathrm{a}'_{6}&=\mathrm{a}_{6}
\end{aligned}
\tag*{(4.3)}
\]
et \uline{modulo 2}
\[
\mathrm{c}_{4}=1\qquad\mathrm{c}_{6}=1\qquad\Delta=\mathrm{a}_{6}\qquad.
\tag*{(4.4)}
\]
\margin{\handwritten{\figure{\begin{tikzpicture}[x=1em,y=1em]
\draw[thin,gray] (0,1.15) -- (.72,1.12);
\draw[thin,gray] (.25,.85) -- (.25,.22);
\draw[thin,gray] (.88,.46) -- (.94,.57);
\draw[thin,gray] (.88,.23) -- (.93,.31);
\draw[thin,gray] (1.5,.43) .. controls (1.65,.68) and (1.82,.31) .. (2.08,.47);
\end{tikzpicture}}}}

5.\qquad \uline{Cas où} $\mathrm{c}_{4}=0$, $\mathrm{p}=2$ \uline{ou} $\mathrm{p}=3$\tikz[baseline=(mark.base)]{\node[inner sep=0pt] (mark) {.};\handwritten{\draw[overlay,thin,gray] ([xshift=.1em,yshift=.86em]mark.center) .. controls +(-.2em,.5em) and +(-.83em,-.8em) .. ([xshift=-.23em,yshift=-.9em]mark.center) .. controls +(.67em,-.12em) and +(0,-.45em) .. ([xshift=.93em,yshift=.5em]mark.center);}}\margin{\handwritten{\tikz[baseline]{\draw[thin,gray] (0,0) .. controls (.12em,.38em) and (.38em,.27em) .. (.58em,.25em);}}}

a) $\mathrm{p}=3$

En caractéristique $3$, lorsque l'équation est mise sou\tikz[baseline=(ending.base)]{\node[inner sep=0pt] (ending) {s};\node[overlay,inner sep=0pt,anchor=base] at (ending.base) {\deleted{r}};\handwritten{\draw[overlay,thin,gray] ([xshift=-.1em,yshift=.85em]ending.center) .. controls +(-.35em,.32em) and +(-.45em,.82em) .. ([xshift=-.67em,yshift=-.05em]ending.center) .. controls +(.2em,-1em) and +(.8em,-.48em) .. ([xshift=.59em,yshift=-.1em]ending.center) .. controls +(0,.75em) and +(.42em,.34em) .. ([xshift=-.1em,yshift=.55em]ending.center);}} forme
% archive sheet number: 23
